La prima parte dell'esperienza consiste nell'apportare degli aggiustamenti preliminari per preparare l'apparato sperimentale.
Dentro il laboratorio è già presente l'apparato che si può schematicamente osservare nell'immagine (\ref{App}). Per prima cosa bisogna misurare la distanza che il fascio di luce proveniente dalla sorgente laser compie durante il tragitto. Il laser infatti colpisce lo specchio rotante, arriva al primo specchio piano, viene riflesso nel secondo specchio piano ed infine arriva allo specchio concavo. Si può schematizzare questa distanza utilizzando le sorgenti immagini proveniente dalle varie riflessioni. Questa misura è la distanza $D$. Successivamente, con l'ausilio di una squadretta bianca, bisogna verificare che la luce laser sia allineata al centro dello specchio rotante con una tolleranza di $(1-2)\,mm$. \\


Dopodiché si inserisce la prima lente $L_1$ con lunghezza focale $f_1=48\,mm$ a $7 \,cm$ dalla sorgente assicurandosi che il fascio di luce non venga deviato e rimanga allineato con il centro dello specchio rotante. Dopo aver inserito la prima lente si inserisce il porta beam splitter a $18 \,cm$ dalla sorgente. Nel posizionare le lenti e il porta beam splitter bisogna sempre assicurarsi di seguire la guida metallica in modo da allineare tutte le ottiche e non propagare errori sistematici. In seguito si posiziona la lente $L_2$ con una lunghezza focale di $f_2=252\, mm$ ad una distanza di $37.8\,cm$ dalla sorgente. Analogamente al posizionamento della lente $L_1$ bisogna assicurarsi che il fascio di luce non sia stato deviato lungo il suo tragitto e sia ancora centrato nello specchio rotante. In caso contrario è possibile agire sul piano del supporto delle lenti ($L_1$ ed $L_2$) per modificare la direzione del fascio in quanto posizionate magneticamente sul proprio supporto. Una volta posizionate tutte le ottiche lungo la guida magnetica bisogna agire sulla cinghia di trasmissione per rendere la faccia piana dello specchio rotante perpendicolare al fascio di luce. In questo modo è possibile indirizzare il fascio di luce verso lo specchio $S_1$ e successivamente verso lo specchio $S_2$. Successivamente bisogna agire sulle viti presenti dietro gli specchi per modificare la loro inclinazione e indirizzare la riflessione del fascio di luce al centro dello specchio concavo.\\

Sulla superficie dello specchio concavo si forma un'immagine della sorgente la cui dimensione può essere modificata spostando la lente $L_2$ lungo la guida magnetica. L'obbiettivo è infatti ridurre al minimo la dimensione dell'immagine presente sullo specchio concavo in modo tale da individuare la posizione dell'immagine all'interno del microscopio con meno ambiguità. 
Per riuscire a misurare la posizione dell'immagine virtuale formata dallo specchio rotante all'interno del microscopio è necessario allineare i cammini di andata e di ritorno del fascio di luce. Con l'ausilio di una foglio di carta traslucida è possibile individuare entrambi i fasci di luce (di andata e di ritorno), ricordandosi che nel momento in cui si copre il fascio di luce di andata quello di ritorno smette di esistere (poiché la luce non raggiunge lo specchio concavo. Effettuando piccoli aggiustamenti alle viti dello specchio concavo e attraverso l'aiuto di un'altra persona che utilizza la carta traslucida per individuare la posizione del fascio di luce di ritorno rispetto al fascio di luce di andata bisogna allineare i due punti. Per farlo, si è verificata la posizione dei due fasci con la carta traslucida in tre diversi posizione lungo il tragitto: vicino allo specchio rotante, in una posizione intermedia tra i due specchi piani e lo specchio concavo ed infine vicino ai due specchi piani.\\

Successivamente, se sono stati eseguiti correttamente tutti gli aggiustamenti preliminari, posizionando la carta traslucida sopra il foro in cui va inserito il microscopio si dovrebbe vedere un punto luminoso. All'interno del microscopio è presente una griglia per agevolare la presa dati, per questo motivo è utile spostare il porta beam splitter fino a visualizzare il punto luminoso al centro del foro, in modo tale da allinearlo con la griglia.\\

Per questioni di sicurezza bisogna inserire sulla guida magnetica un filtro polarizzante tra il laser e la prima lente $L_2$ e successivamente si può montare il microscopio. Dopo aver acceso il motore responsabile della rotazione dello specchio e dopo aver impostato la frequenza di rotazione intorno ai $20-30 \, Hz$ si può togliere il filtro polarizzante e guardare all'interno del microscopio. All'interno si può visualizzare un punto luminoso che con l'aumentare della frequenza di rotazione dello specchio si sposta verso il basso o l'alto in base al senso di rotazione dello specchio.
